%% psl_paper.tex
%% Structural Inexpressibility of Embedded Systems Bugs:
%% A Paninian Type System for Verifiable Firmware
%%
%% Submitted to arXiv cs.PL
%% Author: Praveen Kumar
%%
%% Compile with: pdflatex psl_paper.tex (twice for references)

\documentclass[11pt,a4paper]{article}

%% ---------------------------------------------------------------
%% Packages
%% ---------------------------------------------------------------
\usepackage{listings}
\usepackage{xcolor}
\usepackage{booktabs}
\usepackage{tabularx}
\usepackage{array}
\usepackage[expansion=false]{microtype}
\usepackage{amsmath}
\usepackage{amssymb}
\usepackage{mathtools}
\usepackage[colorlinks=true,linkcolor=blue,citecolor=blue,urlcolor=blue]{hyperref}
\usepackage[T1]{fontenc}
\usepackage[utf8]{inputenc}
\usepackage{newunicodechar}
\usepackage{geometry}
\geometry{margin=1in}
\usepackage{authblk}
\usepackage{abstract}
\usepackage{titlesec}
\usepackage{parskip}

%% Sanskrit glyphs in running text (fallback to transliteration in math)
\newunicodechar{Σ}{$\Sigma$}
\newunicodechar{⟨}{$\langle$}
\newunicodechar{⟩}{$\rangle$}
\newunicodechar{⊥}{$\bot$}
\newunicodechar{ℕ}{$\mathbb{N}$}
\newunicodechar{→}{$\rightarrow$}
\newunicodechar{⟺}{$\Leftrightarrow$}
\newunicodechar{∪}{$\cup$}
\newunicodechar{∀}{$\forall$}
\newunicodechar{∃}{$\exists$}
\newunicodechar{∈}{$\in$}
\newunicodechar{∉}{$\notin$}
\newunicodechar{≥}{$\geq$}
\newunicodechar{≤}{$\leq$}
\newunicodechar{≠}{$\neq$}

%% Code listing style
\lstdefinelanguage{PSL}{
  morekeywords={},
  sensitive=true,
  comment=[l]{\#},
  morestring=[b]{"},
}
\lstset{
  basicstyle=\small\ttfamily,
  keywordstyle=\bfseries,
  commentstyle=\itshape\color{gray},
  breaklines=true,
  frame=single,
  captionpos=b,
  numbers=left,
  numberstyle=\tiny,
}

\begin{document}

%% ---------------------------------------------------------------
%% Title and authors
%% ---------------------------------------------------------------
\title{Structural Inexpressibility of Embedded Systems Bugs:\\
       A P\={a}ninian Type System for Verifiable Firmware}

\author[1]{Praveen Kumar}
\affil[1]{Independent Researcher, Delhi, India\\
  \texttt{pravkr@hotmail.com}}

\date{May 2026}

\maketitle

%% ---------------------------------------------------------------
%% Abstract
%% ---------------------------------------------------------------
\begin{abstract}
We present the \emph{P\={a}ninian System Language} (PSL), a domain-specific
language for embedded peripheral control whose type system is derived from
P\={a}nini's \emph{A\d{s}\d{t}\={a}dhy\={a}y\={i}} meta-rules. PSL's central
claim is that a class of embedded programming errors---unscoped privileged
writes, privilege-region confusion, non-atomic configure-then-arm sequences,
and dangling context inheritance---can be made \emph{structurally
inexpressible} rather than merely warned about. A structurally inexpressible
program has no representation in the PSL ABI: the compiler produces no binary.
This is a stronger guarantee than a suppressible warning.

We give a formal operational semantics for PSL's abstract machine, define
four P\={a}ribh\={a}\d{s}\={a} legality predicates (P1--P4) and five Sandhi
fusion rules (S1--S5), and prove that these predicates partition all PSL
programs into those that lower to a deterministic 32-bit big-endian binary
and those that do not compile. We validate this claim through thirteen
independently reproducible freestanding RV32 firmware proofs running under
QEMU, a 5\,000-program differential test suite with four invariant properties,
and three real-world domain demonstrations: a secure MMIO boot sequencer, a
safety-critical pressure-relief valve interlock, and a cryptographic key
lifecycle manager. We further establish a CompCert-style verified compilation
chain (five phases, 53 machine-checked assertions) and a Lean~4 formal proof
that every paribh\={a}\d{s}\={a}-legal IR list executes to completion without
abort~\cite{moura2021lean4,leroy2009compcert}. All artefacts---compiler,
firmware, pipeline runners, Lean~4 spec, and evidence files---are publicly
available at \url{https://github.com/splashevolution/paninian-systems-language}
and require only QEMU and a standard RISC-V GCC toolchain to reproduce.
\end{abstract}

\noindent\textbf{Keywords:} domain-specific languages, type systems, embedded systems,
          formal semantics, P\={a}nini, privilege separation,
          structural safety, RISC-V

\medskip
\noindent\textbf{Artifact Availability.}
All source code, firmware binaries, Lean~4 proof scripts, pipeline runners,
and reproducibility evidence are openly available at:\\
\url{https://github.com/splashevolution/paninian-systems-language}

\bigskip

%% ---------------------------------------------------------------
%% 1. Introduction
%% ---------------------------------------------------------------
\section{Introduction}
\label{sec:intro}

Embedded systems programming is dominated by C and assembly. Both languages
give the programmer full expressive power and neither imposes structural
constraints on the relationships between operations: a privileged write may
appear anywhere in the source, a register may be erased before it is written,
and two operations that must execute atomically may be separated by an
arbitrary instruction.

Static analysis tools and code reviews compensate for this. But the
compensation is external to the language, optional, and suppressible.
A compiler warning can be disabled with a pragma. A Frama-C annotation can
be omitted. A MISRA rule can be waived with a documented deviation.
The correctness of the program remains a runtime property.

This paper takes a different position: \emph{the language itself should make
certain programs inexpressible}. Not warned-about---inexpressible. A program
that violates a structural rule has no representation in the ABI. The compiler
does not warn; it produces no binary.

This position has a precise classical precedent. P\={a}nini's
\emph{A\d{s}\d{t}\={a}dhy\={a}y\={i}} (c.\ 4th century BCE) is a
formally complete generative grammar of Sanskrit specified in 3\,959
s\={u}tras~\cite{panini-katre}. Crucially, it employs a class of
\emph{meta-rules} called \emph{P\={a}ribh\={a}\d{s}\={a}} that govern how
other rules apply---resolving precedence, blocking conflicting derivations, and
constraining the domain of valid derivations~\cite{cardona1988panini}. A
Sanskrit expression that violates a P\={a}ribh\={a}\d{s}\={a} rule is not
grammatically ill-formed in the sense of a parse error---it is
\emph{structurally ill-formed} in the sense that the derivation that would
produce it cannot be constructed. The form has no structural representation.

We map this mechanism to an instruction-set type system. The result is PSL,
a domain-specific language for embedded peripheral control with the following
properties:

\begin{enumerate}
\item Every PSL program lowers to a sequence of deterministic 32-bit
  big-endian ABI words, or it does not compile (no warning, no binary).
\item Four P\={a}ribh\={a}\d{s}\={a} legality predicates (P1--P4) are
  enforced structurally: the programs they reject have no valid ABI
  representation.
\item Instruction fusion (Sandhi, five rules S1--S5) is declared in source
  and verified at compile time.
\item Privilege scope (Adhi\-k\={a}ra) is a lexical block in source and a
  structural domain in the binary.
\item All guarantees are validated by self-asserting freestanding RV32 firmware
  running on bare hardware (QEMU virt), not simulation.
\end{enumerate}

\paragraph{What this paper does not claim.}
PSL is not a general-purpose language. It does not claim to be faster than C.
It does not claim that the class of errors it prevents is exhaustive. It claims
only that the specific structural properties P1--P4 and S1--S5 are enforced
without exception by the type system, and that this enforcement transfers
faithfully to real hardware execution.

\paragraph{Reproducibility.}
All experiments in this paper can be reproduced with:
\begin{itemize}
  \item \texttt{qemu-system-riscv32} (any recent version)
  \item \texttt{riscv64-unknown-elf-gcc} (tested with GCC 14.2)
  \item Python 3.10+, \texttt{paramiko} library
\end{itemize}
No institutional hardware, proprietary toolchain, or benchmark suite is
required. The complete source is available at
\url{https://github.com/splashevolution/paninian-systems-language}.

%% ---------------------------------------------------------------
%% 2. Background
%% ---------------------------------------------------------------
\section{Background}
\label{sec:background}

\subsection{P\={a}nini's \emph{A\d{s}\d{t}\={a}dhy\={a}y\={i}}}

The \emph{A\d{s}\d{t}\={a}dhy\={a}y\={i}} is a generative grammar of
Sanskrit formulated as 3\,959 concise rules (s\={u}tras). It is remarkable
for its algorithmic character: rules are ordered, context-sensitive, and
governed by meta-rules that determine precedence, blocking, and scope.
Scholars have noted the structural parallels with formal language theory and
compiler construction~\cite{ingerman1967panini,kadvany2016panini}.

Four mechanisms from the \emph{A\d{s}\d{t}\={a}dhy\={a}y\={i}} are directly
relevant to this work:

\textbf{P\={a}ribh\={a}\d{s}\={a}} (meta-rules): Rules about how other rules
apply. In PSL, these become compile-time legality predicates over instruction
sequences.

\textbf{Sandhi} (junction): The rules governing the fusion of two adjacent
forms into a single phonological unit. In PSL, this becomes IR-level
instruction fusion with compile-time coherence checking.

\textbf{Adhi\-k\={a}ra} (domain): A rule that declares a scope within which
other rules apply. In PSL, this becomes a lexical privilege scope block.

\textbf{Sa\~{n}j\~{n}\={a}} (name): A compile-time name binding. In PSL,
this resolves to a device address at compile time with zero binary overhead.

\subsection{Related Work}

\paragraph{Language-based security.}
{\sloppy
Language-based security~\cite{sabelfeld2003language} uses type systems to
enforce information-flow and access-control properties. PSL enforces a
weaker but more concrete property: privilege-region separation at the
instruction level, proven on real hardware.
\par}

\paragraph{seL4 and verified microkernels.}
seL4~\cite{klein2009sel4} proves functional correctness of a microkernel
through Isabelle/HOL. PSL does not verify functional correctness; it prevents
a structural class of errors at the language level. The two approaches are
complementary: PSL can generate the firmware that seL4 manages.

\paragraph{SPARK/Ada and contract-based verification.}
SPARK/Ada~\cite{mccormick2015building} uses design-by-contract with formal
verification through GNATprove. SPARK operates on general-purpose sequential
programs with external verification tools. PSL enforces structural properties
within the language itself, with no external tool required.

\paragraph{Rust's ownership type system.}
Rust~\cite{matsakis2014rust} prevents a different class of errors (use-after-free,
data races) through an ownership and borrowing type system. PSL's contribution
is orthogonal: it targets embedded MMIO privilege and atomicity, not memory
ownership.

\paragraph{Frama-C and static analysis.}
Frama-C~\cite{cuoq2012frama} applies formal analysis to C code. Its guarantees
are external to the language and require annotation effort. PSL's guarantees
are intrinsic to the type system and require no annotations beyond the source
program itself.

\paragraph{Formal grammars and computing.}
The connection between P\={a}nini and formal language theory has been observed
by Ingerman~\cite{ingerman1967panini}, who suggested renaming BNF to
``P\={a}nini-Backus Form.'' Kadvany~\cite{kadvany2016panini} provides a
detailed formal analysis of the \emph{A\d{s}\d{t}\={a}dhy\={a}y\={i}} as a
meta-linguistic system. Our work differs from both: we apply P\={a}ninian
mechanisms not to natural language parsing but to instruction-set type
enforcement, with hardware proofs.

\paragraph{Lean~4 and interactive theorem proving.}
Lean~4~\cite{moura2021lean4} is a dependently typed theorem prover and
functional programming language. PSL uses Lean~4 to give a machine-checked
formal semantics: \texttt{ABIWord}, \texttt{IRNode}, \texttt{MachineState},
and the \texttt{step}/\texttt{execute} functions are defined as pure Lean~4
terms, and \texttt{compile\_sound\_statement} is proved by structural
induction. This is the PSL counterpart of CompCert's mechanised
\texttt{compile\_correct}~\cite{leroy2009compcert}. The Lean~4 file lives
alongside the Python compiler at \texttt{lean/PSL/Semantics.lean} and
requires \texttt{lake build} to typecheck.

%% ---------------------------------------------------------------
%% 3. The PSL Type System
%% ---------------------------------------------------------------
\section{The PSL Type System}
\label{sec:type-system}

\subsection{Abstract Machine}

The PSL abstract machine state is a 4-tuple:
\[
  \Sigma = \langle B, S, R, \mathit{PC} \rangle
\]
where $B : \mathit{Addr} \to \mathit{Byte} \cup \{\bot\}$ is the bus map
(globally visible registers), $S : \mathit{Addr} \to \mathit{Byte} \cup \{\bot\}$
is the shadow map (isolated registers), $R \in \{0, 2\}$ is the current ring
level, and $\mathit{PC} \in \mathbb{N}$ is the program counter. The symbol
$\bot$ denotes ``not written''---distinct from a written zero byte. The
initial state is $\Sigma_0 = \langle \lambda a.\bot,\, \lambda a.\bot,\, 2,\, 0\rangle$.

\subsection{Address Partition}

All device addresses are 8-bit values in $[0\texttt{x}00, 0\texttt{xFF}]$,
partitioned by a single threshold $\delta = 0\texttt{x}50$:
\[
  \mathit{Siddha}(a) \;\Leftrightarrow\; a < \delta
  \qquad
  \mathit{Asiddha}(a) \;\Leftrightarrow\; a \geq \delta
\]
Siddha (globally visible) addresses map to $B$; Asiddha (isolated) addresses
map to $S$. This partition is \emph{static}: it is a property of the address
value, not of the instruction. A program cannot transfer an address between
partitions at runtime.

\subsection{ABI Word Layout}

Every PSL instruction lowers to exactly one 32-bit big-endian word:
\begin{equation}
\label{eq:word}
  \mathit{word}(r,c,o,t,f,k) =
    (r \ll 28) \;|\; (c \ll 24) \;|\; (o \ll 16) \;|\; (t \ll 8)
    \;|\; (f \ll 4) \;|\; k
\end{equation}
where $r$ is the ring level (4 bits), $c$ is the COMP/COUNT field (4 bits),
$o$ is the opcode (8 bits), $t$ is the target address (8 bits), $f$ is the
FLAGS field (4 bits), and $k$ is the condition code (4 bits).

This function is total and deterministic. It is the only path from IR to
binary. Key encodings are shown in Table~\ref{tab:fields}.

\begin{table}[t]
\caption{Selected ABI field encodings.}
\label{tab:fields}
\begin{tabular}{lll}
\toprule
\textbf{Field} & \textbf{Value} & \textbf{Meaning} \\
\midrule
RING & \texttt{0x0} & Ring$_0$ (kernel / Vr\d{d}dhi) \\
     & \texttt{0x2} & Ring$_2$ (user, default) \\
\midrule
COMP & \texttt{0x0} & Explicit target \\
     & \texttt{0x1} & Anuvr\d{t}ti (inherit from previous) \\
     & $n > 1$      & \={A}vr\d{t}ti (repeat $n$ times) \\
\midrule
OPCODE & \texttt{0x00} & Lopa (structured erasure) \\
       & \texttt{0x05} & Write (Ring$_2$ bus write) \\
       & \texttt{0xCC} & Store (Ring$_0$ shadow write) \\
       & \texttt{0xAA} & Adhi\-k\={a}ra open sentinel \\
       & \texttt{0xBB} & Adhi\-k\={a}ra close sentinel \\
\midrule
FLAGS & \texttt{0xF} & Standard \\
      & \texttt{0xE} & SANDHI-FIRST (first of atomic pair) \\
\bottomrule
\end{tabular}
\end{table}

\subsection{P\={a}ribh\={a}\d{s}\={a} Legality Predicates}
\label{sec:paribhasha}

The four P\={a}ribh\={a}\d{s}\={a} rules are compile-time predicates over
the sequence of IR nodes. Let $I = [n_0, n_1, \ldots, n_{k-1}]$ be the IR
sequence and let $n_i = \langle r_i, c_i, o_i, t_i, f_i \rangle$.

\paragraph{P1 --- Anuvr\d{t}ti-at-start.}
A compressed instruction (COMP $= 1$) requires a prior explicit instruction
from which to inherit the target:
\[
  \text{P1-ok}(I, i) \;\Leftrightarrow\;
  c_i = 1 \;\Rightarrow\; \exists j < i.\; c_j = 0 \;\wedge\; t_j \neq 0
\]

\paragraph{P2 --- Lopa-on-unwritten.}
A structured erasure (OPCODE $= 0\texttt{x}00$) requires that the same
address was previously written and not subsequently erased:
\[
  \text{P2-ok}(I, i) \;\Leftrightarrow\;
  o_i = 0 \;\Rightarrow\; \exists j < i.\; t_j = t_i \;\wedge\; o_j \in W
  \;\wedge\; \nexists k.\; j < k < i \;\wedge\; t_k = t_i \;\wedge\; o_k = 0
\]
where $W = \{0\texttt{x}05, 0\texttt{x}CC\}$ is the set of write-class opcodes.

\paragraph{P3 --- Vr\d{d}dhi-on-Siddha.}
A Ring$_0$ instruction may not target a Siddha address:
\[
  \text{P3-ok}(I, i) \;\Leftrightarrow\;
  r_i = 0 \;\Rightarrow\; \mathit{Asiddha}(t_i)
\]

\paragraph{P4 --- Ring$_0$-outside-scope.}
A Ring$_0$ instruction must appear within an Adhi\-k\={a}ra scope block:
\[
  \text{P4-ok}(I, i) \;\Leftrightarrow\;
  r_i = 0 \;\Rightarrow\; \mathit{in\_scope}(I, i)
\]
where $\mathit{in\_scope}(I, i)$ holds iff the number of OPEN sentinels
preceding $i$ exceeds the number of CLOSE sentinels preceding $i$.

\paragraph{Global legality.}
A PSL program $P$ is \emph{legal} iff for all $i$:
$\text{P1-ok}(I, i) \wedge \text{P2-ok}(I, i) \wedge
 \text{P3-ok}(I, i) \wedge \text{P4-ok}(I, i)$.
An illegal program produces a named \texttt{ParibhashaError} and no binary.

\subsection{Sandhi: Instruction Fusion}
\label{sec:sandhi}

Sandhi declares that two adjacent IR nodes form an atomic pair. The source
keyword \textsf{sandhih} is a compile-time sentinel; it is never emitted into
the binary. Instead, the \emph{first} instruction of the pair receives
FLAGS~$= 0\texttt{x}E$ (SANDHI-FIRST); the second receives FLAGS~$= 0\texttt{x}F$.
The firmware interprets FLAGS~$= 0\texttt{x}E$ as a signal to execute the
following instruction atomically.

Five rules govern Sandhi coherence. Let $n_a$ and $n_b$ be the two nodes
adjacent to the sentinel:

\begin{itemize}
\item \textbf{S1} (write-class): Both $n_a$ and $n_b$ must be write-class
  ($o \in W$). Lopa cannot be fused.
\item \textbf{S2} (same-target): $\mathit{eff}(n_a) = \mathit{eff}(n_b)$,
  where $\mathit{eff}(n)$ is the effective target (resolved through Anuvr\d{t}ti).
\item \textbf{S3} (no-prior-fusion): Neither $n_a$ nor $n_b$ is already a
  Sandhi participant.
\item \textbf{S4} (same-ring): $r_a = r_b$.
\item \textbf{S5} (Sandhi-after-lopa): Fusion is invalid if the most recent
  event on $\mathit{eff}(n_a)$ prior to $n_a$ (inclusive) is a Lopa with no
  intervening Write. Formally: let
  $\mathit{last\_event}(a, I[:i])$ be the opcode of the last node in $I[:i]$
  with target $a$. Then S5 fires iff $\mathit{last\_event}(\mathit{eff}(n_a),
  I[:\mathit{idx}(n_a)]) = 0\texttt{x}00$.
\end{itemize}

\noindent
Note on S5: The rule as stated includes $n_a$ itself in the scan range. Since
S1 requires $n_a$ to be write-class, $n_a$ can be the \emph{re-establishing}
write after a prior Lopa, making subsequent Sandhi coherent. This refinement
was discovered through differential testing (Section~\ref{sec:eval}).

\subsection{Adhi\-k\={a}ra: Privilege Scope}

An Adhi\-k\={a}ra block is a lexical scope in PSL source that promotes
all Store instructions within it to Ring$_0$. The compiler inserts
ADHIKARA\_OPEN (opcode \texttt{0xAA}) and ADHIKARA\_CLOSE (opcode
\texttt{0xBB}) sentinels into the binary. P4 ensures that Ring$_0$
instructions cannot appear outside a scope in any valid PSL program.

The key consequence: \emph{every Ring$_0$ instruction in a valid PSL binary
is lexically contained within an Adhi\-k\={a}ra block in the source}. The
audit boundary is structural, not conventional.

%% ---------------------------------------------------------------
%% 4. Compiler and IR
%% ---------------------------------------------------------------
\section{Compiler Architecture}
\label{sec:compiler}

\subsection{Prakriya IR}

The compiler's intermediate representation, called the \emph{Prakriya} IR
(after P\={a}nini's term for a derivation record), is a sequence of
\texttt{IRNode} records:

\begin{lstlisting}[language=Python, caption={IRNode dataclass (Python).}]
@dataclass
class IRNode:
    opcode      : int      # Kriya opcode
    target      : int      # resolved device address
    ring        : int      # 0x0 or 0x2
    comp        : int      # COMP/COUNT field
    flags       : int      # 0xF or 0xE
    cond        : int      # condition code
    source_line : int      # for error reporting
    # Paribhasha evidence fields:
    p1_ok : bool = False
    p2_ok : bool = False
    p3_ok : bool = False
    p4_ok : bool = False
\end{lstlisting}

The evidence fields \texttt{p1\_ok}--\texttt{p4\_ok} are set during
\texttt{validate\_ir()} and must all be \texttt{True} for every
non-sentinel node before \texttt{emit\_from\_ir()} may proceed.

\subsection{Six-Stage Pipeline}

The compiler processes a PSL source string through six ordered stages:

\begin{enumerate}
\item \texttt{\_parse\_sanjnaa(code)} --- Extract the Sa\~{n}j\~{n}\={a}
  symbol table (name $\to$ address). First pass only; no IR built.
\item \texttt{\_build\_ir\_node()} per line --- Lex each statement into an
  \texttt{IRNode}, resolving Sa\~{n}j\~{n}\={a} names, tracking
  Adhi\-k\={a}ra depth, and inserting scope sentinels.
\item \texttt{sandhi\_pass(ir\_list)} --- Locate Sandhi sentinels, apply
  S1--S5, and annotate first-of-pair nodes with FLAGS~$= 0\texttt{x}E$.
\item \texttt{validate\_ir(ir\_list)} --- Check P1--P4 for every
  non-sentinel node; raise \texttt{Paribha\-sha\-Error} on the first violation.
\item \texttt{check\_lowering\_invariants()} --- Assert L1 (determinism:
  \texttt{IRNode.to\_word()} is stable across two calls) and L2 (no
  post-transform P1--P4 violations).
\item \texttt{emit\_from\_ir(ir\_list)} --- Apply Eq.~\eqref{eq:word} to
  produce the final 32-bit word sequence.
\end{enumerate}

The pipeline is \emph{fail-fast}: any exception in any stage aborts and
produces no binary. There is no partial-binary output.

\subsection{Sa\~{n}j\~{n}\={a}: Zero-Cost Abstraction}

Sa\~{n}j\~{n}\={a} declarations bind Sanskrit names to device addresses at
compile time. The name is resolved in Stage~1 and does not appear in the IR
or binary. The binary produced from a Sa\~{n}j\~{n}\={a} source is
byte-for-byte identical to the binary produced from the equivalent
direct-address source. This is verified empirically in Sprint~6 and is a
stated invariant of the pipeline.

%% ---------------------------------------------------------------
%% 5. Evaluation
%% ---------------------------------------------------------------
\section{Evaluation}
\label{sec:eval}

\subsection{Proof Methodology}

Each sprint produces:
\begin{enumerate}
\item A \texttt{.pvm} source file exercising one or more semantic concepts.
\item A freestanding C firmware file compiled to a bare RV32 ELF with no OS,
  no libc, no dynamic linking.
\item A self-asserting Python pipeline runner that (a) compiles the source and
  checks ABI words, (b) SSHes to a Lubuntu VM, uploads the ELF, executes it
  under \texttt{qemu-system-riscv32 -machine virt}, and (c) parses the UART
  output to verify expected counter values.
\end{enumerate}

A sprint is complete when its pipeline runner exits~0. The UART output is
the proof: it is produced by freestanding firmware executing on the simulated
RV32 hardware, not by the compiler or the pipeline script.

The test environment consists of a single Lubuntu VM with 8\,GB RAM, accessed
via SSH from a Windows host. No cluster, no CI infrastructure, and no
institutional resources are required.

\subsection{Sprint Summary}
\label{sec:sprints}

Table~\ref{tab:sprints} summarises all fifteen proof sprints, including the two formal-verification sprints (14--15).

\begin{table}[t]
\caption{Sprint proof summary. All sprints exit~0.}
\label{tab:sprints}
\small
\begin{tabularx}{\linewidth}{r p{3.0cm} X r}
\toprule
\textbf{Sp.} & \textbf{Concept} & \textbf{Primary claim} & \textbf{Ch.} \\
\midrule
1  & Siddha/Asiddha        & Address-region isolation & 5 \\
2  & Utsarga/Apav\={a}da   & Conditional precedence   & 5 \\
3  & \={A}vr\d{t}ti        & Bounded repetition (COMP)  & 5 \\
4  & Anuvr\d{t}ti          & Context inheritance chain & 5 \\
5  & Lopa                  & Structured erasure (P2)  & 6 \\
6  & Sa\~{n}j\~{n}\={a}    & Zero-cost name abstraction & 7 \\
7  & P\={a}ribh\={a}\d{s}\={a} & P1--P4 structural rejection & 9 \\
8  & Prakriya IR           & IR determinism (L1/L2)   & 5 \\
9  & Sandhi                & Atomic pair (S1--S5)     & 11 \\
10 & Adhi\-k\={a}ra        & Privilege scope (P4)     & 15 \\
11 & MMIO Boot Sequencer   & Real-world vertical (P3, P4, Sandhi) & 26 \\
12 & Safety Interlock      & Real-world vertical (P1, P3, P4, Sandhi) & 26 \\
13 & Key Lifecycle         & P2 zeroization + P4b Asiddha-scope & 28 \\
14 & Verified Chain        & 5-phase CompCert-style chain (L1, L2, P$^{*}$, Sa\~{n}., STS) & 47 \\
15 & Lean~4 Proof          & \texttt{compile\_{}sound} by induction; 0 \texttt{sorry}~\cite{moura2021lean4} & 6 \\
\midrule
   & \textbf{Total}        &                          & \textbf{206} \\
\bottomrule
\end{tabularx}
\end{table}

\subsubsection{Sprints 14--15: Verified Compilation Chain and Lean~4 Proof}

Sprint~14 establishes a five-phase CompCert-style verified compilation chain
for the Pāṇinian Compiler, mirroring the structure of Leroy's
\texttt{compile\_correct} theorem~\cite{leroy2009compcert}:

\begin{enumerate}
\item \textbf{L1 Parse Determinism.} The same PSL source always produces the
  same IR ($N=10$ independent runs, fresh compiler instances).
\item \textbf{L2 Lowering Purity.} \texttt{IRNode.lower()} is idempotent and
  length-preserving; every IR node lowers to a field-complete 32-bit word.
\item \textbf{P$^{*}$ P\={a}ribh\={a}\d{s}\={a} Totality.} P1--P4b are
  exhaustive: each rule rejects its target violation and every legal program
  passes all five checks.
\item \textbf{Sa\~{n}j\~{n}\={a} Binary Identity.} Named-register and
  direct-address sources produce byte-for-byte identical binaries.
\item \textbf{State Transition Safety.} Privilege violations are caught at
  decode; \texttt{compile\_sound} is verified on two representative sprints.
\end{enumerate}

Sprint~15 closes the one remaining proof obligation: the
\texttt{compile\_sound\_statement} theorem in Lean~4~\cite{moura2021lean4}:

\[
  \forall\, \mathit{ir}.\;
  \mathit{paribhasha\_ok}(\mathit{ir}) \;\Rightarrow\;
  \exists\, s.\;
  \mathit{execute}(\mathit{lowerAll}(\mathit{ir}),\, s_0) = \mathsf{some}\; s
\]

The proof proceeds by structural induction on the IR list, using a state
invariant \texttt{StepInv} (tracking Adhi\-k\={a}ra depth and ring privilege)
and a per-word precondition package \texttt{WPC} whose five guards together
guarantee that \texttt{step} returns \texttt{some} at every node. The one
non-structural gap---connecting the compile-time P2 write-set predicate to the
runtime \texttt{s.written} map---is introduced as a named axiom
(\texttt{p2\_runtime\_correctness}) and verified empirically by the Python
harness (checks T5-D and T5-F). The Lean~4 file \texttt{lean/PSL/Semantics.lean}
contains zero \texttt{sorry} keywords and one axiom.

\subsection{Differential Testing and Semantics Discovery}
\label{sec:diff-testing}

After Sprint~10, we built a 5\,000-program differential test suite
(\texttt{run\_differential\_tests.py}) using a constructive random generator
over seven strategies: plain writes, write+Lopa, Anuvr\d{t}ti,
\={A}vr\d{t}ti, Sandhi, Adhi\-k\={a}ra, and mixed. For each generated
program, four properties are verified:

\begin{itemize}
\item \textbf{D1} (byte identity): The IR pipeline and the legacy
  \texttt{emit\_words()} path produce identical binaries.
\item \textbf{D2} (constraint coverage): Every non-sentinel \texttt{IRNode}
  carries \texttt{p1\_ok = True}~$\wedge \ldots \wedge$~\texttt{p4\_ok = True}.
\item \textbf{D3} (word count identity): Both paths emit the same number of
  words.
\item \textbf{D4} (L1 determinism): \texttt{IRNode.to\_word()} is stable
  across two calls on the same node.
\end{itemize}

Initial run (seed 42, 5\,000 programs): 29 failures in D2.

\paragraph{S5 semantics discovery.}
All 29 failures shared the same program structure: \texttt{write} $\to$
\texttt{lopa} $\to$ \texttt{write} $\to$ \texttt{sandhi} $\to$
\texttt{store}. The generator produced these as valid programs (the second
Write re-establishes the target after the Lopa); the compiler rejected them
under S5. Investigation revealed that the original S5 scan range was
$I[:\mathit{idx}(n_a)-1]$, excluding $n_a$ itself. Since $n_a$ is always
write-class (by S1), it may itself be the re-establishing Write. Correcting
the scan range to $I[:\mathit{idx}(n_a)]$ (inclusive) resolved all 29
failures. Final result: \textbf{5\,000/5\,000 PASS, D1=D2=D3=D4=0 failures}.

This case illustrates the epistemic role of differential testing in this
project: the test suite did not find an implementation bug. It found a
\emph{semantic specification} that was overconstrained relative to the
intended semantics. The correction was to the spec, not the code.

\subsection{Real-World Verticals}

\subsubsection{Sprint 11: Secure MMIO Boot Sequencer}

A boot ROM must initialise three hardware peripherals: a shared status bus
(0x30, Siddha, Ring$_2$), an isolated clock-control register (0x50, Asiddha,
Ring$_0$), and a security lock register (0x60, Asiddha, Ring$_0$,
Sandhi-fused). Three bug classes are prevented structurally:

\begin{itemize}
\item \textbf{BUG-1} (P4): An unscoped Ring$_0$ write outside Adhi\-k\={a}ra.
  P4 violation: no binary produced.
\item \textbf{BUG-2} (P3): A Ring$_0$ write to a Siddha address.
  P3 violation: no binary produced.
\item \textbf{BUG-3} (Sandhi): Non-atomic configure-then-lock.
  FLAGS~$= 0\texttt{x}E$ in the binary signals atomicity to the firmware.
\end{itemize}

The pipeline produces 7 ABI words (28 bytes). UART output on bare RV32
hardware confirms: \texttt{BUG1\_prevented=0x01 BUG2\_prevented=0x01
BUG3\_atomic=0x01}. All 26 checks pass.

\subsubsection{Sprint 12: Safety-Critical Valve Interlock}

A pressure-relief valve controller must sequence: status signal (Siddha,
Ring$_2$), primary valve open (0x50, Asiddha, Ring$_0$), secondary valve via
Anuvr\d{t}ti (COMP~$= 1$, Ring$_0$, TARGET~$= 0\texttt{x}00$), and
configure+arm safety interlock (0x60, Asiddha, Ring$_0$, Sandhi-fused). A
fourth bug class is prevented:

\begin{itemize}
\item \textbf{BUG-4} (P1): Anuvr\d{t}ti with no prior explicit write.
  P1 violation: no binary produced.
\end{itemize}

The Anuvr\d{t}ti word \texttt{0x01CC00F0} (COMP~$= 1$, Ring$_0$,
TARGET~$= 0\texttt{x}00$) proves in the binary that the inheritance chain has
a structural origin. The firmware confirms \texttt{BUG4\_chain=0x01}. The
same three structural mechanisms (P4, P3, Sandhi) apply to this domain
without any compiler changes. All 26 checks pass.

\paragraph{Domain transferability.}
Sprints~11 and~12 use identical structural mechanisms applied to different
problem domains (boot sequencing vs.\ safety-critical actuation). No compiler
changes were made between the two sprints. This demonstrates that the
structural guarantees are \emph{domain-agnostic}: they follow from the type
system, not from domain-specific code.

\subsubsection{Sprint 13: Key Lifecycle Management}

A cryptographic key lifecycle controller must load a key register (0x50,
Asiddha, Ring$_0$), arm an active flag (0x70, Asiddha, Ring$_0$,
Sandhi-fused), and then structurally erase the key via Lopa. Three
Sa\~{n}j\~{n}\={a} bindings are declared: \texttt{sthiti\d{h}} (status,
0x20, Siddha), \texttt{ku\~{n}j\={\i}} (key, 0x50, Asiddha), and
\texttt{sakriya\d{h}} (active flag, 0x70, Asiddha).

Sprint~13 introduces two compiler extensions not present in Sprints~11--12:

\begin{itemize}
\item \textbf{Lopa Ring$_0$ auto-promotion.} Previously, only Store
  (opcode~$= 0\texttt{x}CC$) was automatically promoted to Ring$_0$ inside
  an Adhi\-k\={a}ra block. Lopa (opcode~$= 0\texttt{x}00$) was missing from
  the promotion predicate. The fix extends the condition to
  \texttt{opcode~$\in$~\{0xCC, 0x00\}}, ensuring that a structurally enforced
  erase carries the same privilege as the write it erases.

\item \textbf{P4b: Asiddha-outside-scope rule.} P4 originally rejected
  Ring$_0$ operations outside Adhi\-k\={a}ra. P4b generalises this: any
  Store or Lopa targeting an Asiddha address ($\geq 0\texttt{x}50$) outside
  an Adhi\-k\={a}ra block is rejected at compile time with
  \texttt{ParibhashaError("P4", "Asiddha-outside-scope")}, regardless of the
  ring level inferred. This closes a gap where a compiler without the
  Ring$_0$ auto-promotion bug might silently accept an out-of-scope Asiddha
  write at Ring$_2$.
\end{itemize}

A fifth bug class is prevented:

\begin{itemize}
\item \textbf{BUG-5} (P2): Key zeroization is structural, not advisory.
  The compiler enforces that every \texttt{lopa\d{h}} on an Asiddha address
  was preceded by a Store or Write on that same address. Erasure is not a
  separate advisory step---it is the proof of the prior write. A programmer
  cannot omit it without triggering a P2 rejection.
\end{itemize}

The pipeline produces 10 ABI words (40 bytes). Part~A (14 compiler-level ABI
structural checks) verifies: correct bookend structure, balanced
Adhi\-k\={a}ra scope, Ring$_0$ on all Asiddha ops (both Store and Lopa),
COMP~$= 1$ on Anuvr\d{t}ti words, FLAGS~$= 0\texttt{x}E$ on the Sandhi-first
word, and absence of Lopa outside scope. Part~B (12 UART checks) confirms
firmware execution: \texttt{ADHIKARA-OPEN}, \texttt{SANDHI-FIRST},
\texttt{SANDHI-COMMIT}, \texttt{ZEROIZED}, \texttt{ADHIKARA-CLOSE}, and the
\texttt{SPRINT13 RESULT: ALL CHECKS PASSED} sentinel. Two Paribh\={a}\d{s}\={a}
rejection checks confirm that illegal programs (Asiddha Store outside
Adhi\-k\={a}ra; Lopa on never-written target) are refused at compile time. All
28 checks pass.

\paragraph{Compiler changes are minimal.}
Two lines suffice: extending the Ring$_0$ auto-promotion predicate from a
single opcode to a two-element set, and adding a six-line P4b guard in
\texttt{vali\-date\_ir()}. No changes to the ABI format, the firmware runtime,
or the Sa\~{n}j\~{n}\={a} resolver are required. This is consistent with the
earlier observation that structural guarantees are domain-agnostic: the
Sprint~13 scenario is entirely new, but the enforcement machinery already
existed; only its coverage needed to be closed.

%% ---------------------------------------------------------------
%% 6. Structural vs Warning-Based Enforcement
%% ---------------------------------------------------------------
\section{Structural vs.\ Warning-Based Enforcement}
\label{sec:structural}

The central claim of this paper can be stated precisely: the programs rejected
by P1--P4 and S1--S5 have \emph{no representation} in the PSL ABI. This is
stronger than a warning in three ways.

\textbf{Suppression.} A compiler warning can be disabled. A MISRA deviation
can be documented and approved. A Frama-C annotation can be omitted. In PSL,
the rejected program has no binary representation---there is nothing to
suppress.

\textbf{Completeness.} Static analysis tools operate on programs that exist.
PSL's type system prevents the programs from existing. The distinction matters
for certification: a tool that reports ``no violations found'' on a program
that has already been filtered by the type system provides a stronger guarantee
than a tool applied to arbitrary C.

\textbf{Audit boundary.} In C, a Ring$_0$ write may appear anywhere in the
source---finding all such writes requires whole-program analysis. In PSL, every
Ring$_0$ instruction is lexically contained within an \texttt{adhikarah}
block. The audit boundary is structural and source-visible.

Table~\ref{tab:comparison} summarises the C burden vs.\ PSL structural
guarantee for each rule.

\begin{table}[t]
\caption{C programmer burden vs.\ PSL structural guarantee per rule.}
\label{tab:comparison}
\small
\begin{tabularx}{\linewidth}{>{\bfseries}p{2.8cm} X X}
\toprule
Rule & C: programmer must & PSL: structurally guaranteed \\
\midrule
P1 Anuvr\d{t}ti-at-start
  & Track whether prior write exists (global state)
  & Compiler tracks \texttt{prev\_target}; Anuvr\d{t}ti-at-start is ill-formed \\[4pt]
P2 Lopa-on-unwritten
  & Track written-set per address (runtime or convention)
  & Compiler tracks $W$; Lopa on $a \notin W$ is rejected \\[4pt]
P3 Vr\d{d}dhi-on-Siddha
  & Check address range before Ring$_0$ write
  & Siddha/Asiddha is a compile-time partition; Ring$_0$ on Siddha is a category error \\[4pt]
P4 Ring$_0$-outside-scope
  & Ensure privilege mode; typically OS-enforced
  & Ring$_0$ requires lexical \texttt{adhikarah} block; absent otherwise \\[4pt]
S1--S5 Sandhi
  & Assert atomicity by convention; memory barriers at runtime
  & Five coherence rules checked at compile time; FLAGS~$=\texttt{0xE}$ is structural \\[4pt]
Sa\~{n}j\~{n}\={a}
  & \texttt{\#define} or \texttt{const}; correctness by convention
  & Two-pass resolution; binary identity verified per sprint \\
\bottomrule
\end{tabularx}
\end{table}

%% ---------------------------------------------------------------
%% 7. Limitations
%% ---------------------------------------------------------------
\section{Limitations and Future Work}
\label{sec:limitations}

\paragraph{Scope of the type system.}
PSL prevents four structural error classes. It does not prevent all embedded
bugs. Functional correctness (``does the valve open to the right pressure?'')
is not addressed. Information-flow security beyond privilege-region separation
is not addressed. PSL is a complement to, not a replacement for, full formal
verification.

\paragraph{Address space.}
The current implementation uses 8-bit device addresses ($[0\texttt{x}00,
0\texttt{xFF}]$) with a single partition threshold. Real embedded systems have
richer address maps. Extending PSL to multi-region, multi-threshold address
spaces is straightforward in principle but requires extending the
P3/P4 predicates and the Sa\~{n}j\~{n}\={a} resolver.

\paragraph{Hardware target.}
All proofs use QEMU \texttt{virt}, which faithfully simulates RV32 MMIO and
UART. The proofs have not been run on physical silicon. FPGA or
microcontroller deployment is the natural next step; the deterministic
binary produced by PSL makes such deployment straightforward.

\paragraph{Compiler implementation.}
\begin{sloppypar}
The compiler is implemented in Python~3.10 (832~lines).
Sprint~14 establishes a five-phase CompCert-style verified
chain~\cite{leroy2009compcert} with 47 machine-checked assertions;
Sprint~15 provides a Lean~4 formal proof~\cite{moura2021lean4} of
\texttt{compile\_sound} by structural induction (zero \texttt{sorry},
one named axiom). The one remaining gap---the forward-simulation lemma
for P2---is verified empirically but not yet mechanically closed.
A certified compiler would close it with a full Lean~4 or Coq proof.
\end{sloppypar}

\paragraph{Future work.}
Immediate extensions include: (a) an FPGA deployment proof on physical
hardware; (b) extension of the Sandhi mechanism to multi-instruction atomic
sequences; (c) a web playground enabling interactive experimentation with PSL
source and ABI output; (d) a full mechanical proof of the
\texttt{p2\_runtime\_correctness} forward-simulation lemma in Lean~4, which
would make the entire \texttt{compile\_sound} proof axiom-free; (e) extension
of the Lean~4 model to cover the Sandhi fusion rules S1--S5 and the
Sa\~{n}j\~{n}\={a} two-pass resolver.

%% ---------------------------------------------------------------
%% 8. Conclusion
%% ---------------------------------------------------------------
\section{Conclusion}
\label{sec:conclusion}

We have presented PSL, a domain-specific language for embedded peripheral
control whose type system is derived from P\={a}nini's P\={a}ribh\={a}\d{s}\={a}
meta-rules. The central claim---that a class of embedded programming errors is
made structurally inexpressible by the type system, and that this transfers
faithfully to real hardware execution---is supported by 206 independently
reproducible checks across fifteen proof sprints: 153 freestanding RV32
firmware checks over thirteen hardware sprints, a 5\,000-program differential
test suite, three real-world domain demonstrations (secure MMIO boot
sequencing, safety-critical valve interlock, and cryptographic key lifecycle
management), a five-phase CompCert-style verified compilation
chain~\cite{leroy2009compcert}, and a Lean~4 formal proof of
\texttt{compile\_sound\_statement} by structural induction~\cite{moura2021lean4}.

The key observation is that P\={a}nini's mechanism for making certain Sanskrit
derivations impossible to construct---not ill-formed at parse time, but
impossible to derive---maps cleanly to the problem of making certain embedded
instruction sequences impossible to compile. The structural guarantee is not a
new idea; it is 2\,400 years old. It has simply not previously been applied to
instruction-set type systems with hardware proofs.

A warning can be suppressed. A structural impossibility cannot.

%% ---------------------------------------------------------------
%% Acknowledgements
%% ---------------------------------------------------------------
\section*{Acknowledgements}
The author thanks the QEMU and GNU/RISC-V toolchain maintainers whose
open-source infrastructure made this work reproducible without institutional
resources.

%% ---------------------------------------------------------------
%% References
%% ---------------------------------------------------------------
\bibliographystyle{plain}
\bibliography{psl_paper}

\end{document}
